\subsubsection{Requisitos del Módulo de Gestión de Inventario}
\footnotesize
\begin{longtable}{>{\raggedright\arraybackslash}p{0.12\textwidth} >{\raggedright\arraybackslash}p{0.50\textwidth} >{\raggedright\arraybackslash}p{0.18\textwidth} >{\raggedright\arraybackslash}p{0.12\textwidth}}
\toprule
\rowcolor{gray!15}\textbf{ID} & \textbf{Descripción} & \textbf{Categoría} & \textbf{Prioridad} \\ \midrule
\endhead
FR-116 & El sistema permitirá seleccionar una Ubicación Funcional (FL) del árbol jerárquico para vincular el nuevo activo de nivel 6 según la taxonomía ISO 14224. & Functional Suitability & Alta \\
FR-117 & El sistema validará la unicidad del identificador industrial (Tag) a nivel global, impidiendo duplicidades en el catálogo maestro de equipos. & Functional Suitability & Alta \\
FR-118 & El formulario de registro validará la presencia de datos técnicos críticos: identificación, clase, fabricante, modelo, serial y fechas operativas antes de permitir el guardado. & Functional Suitability & Alta \\
FR-119 & El sistema implementará reglas de coherencia cronológica para asegurar que los hitos operativos del activo mantengan una secuencia lógica y veraz. & Functional Suitability & Alta \\
FR-120 & La ficha del activo presentará una sección técnica consolidada con todos los parámetros capturados, permitiendo la actualización de atributos no críticos mediante procesos auditados. & Functional Suitability & Media \\
FR-125 & El sistema permitirá el registro formal de ingresos de material, capturando el repuesto, la cantidad, la ubicación física y vinculando el movimiento a la auditoría maestra. & Functional Suitability & Alta \\
FR-126 & El sistema mantendrá el cálculo dinámico del stock disponible considerando las cantidades físicas y las reservas pendientes para mantenimiento programado. & Functional Suitability & Alta \\
FR-127 & El sistema gestionará las salidas de material asociadas a órdenes de trabajo, garantizando que el consumo de repuestos sea trazable hasta el activo intervenido. & Functional Suitability & Alta \\
FR-128 & El sistema permitirá la reintegración de materiales no utilizados al inventario maestro, gestionando la reversión de costos en las órdenes de trabajo correspondientes. & Functional Suitability & Alta \\
FR-129 & El sistema emitirá notificaciones automáticas ante condiciones de bajo stock o excesos de inventario según las políticas de reposición configuradas. & Functional Suitability & Alta \\
FR-130 & El sistema proporcionará capacidades de consulta y filtrado sobre el historial de movimientos de inventario para fines de auditoría y conciliación. & Functional Suitability & Media \\
FR-131 & Al completar satisfactoriamente una salida de material, el sistema permitirá generar un comprobante técnico de la operación con los metadatos de la transacción. & Functional Suitability & Media \\
FR-137 & El sistema presentará una lista consolidada de activos en estado preliminar o de revisión, con indicadores visuales de urgencia y tiempo en espera. & Functional Suitability & Alta \\
FR-138 & La interfaz de revisión técnica destacará los campos obligatorios faltantes según la clase de activo, guiando al ingeniero en la completitud de la ficha técnica. & Functional Suitability & Alta \\
FR-139 & El sistema permitirá la edición controlada de los atributos técnicos durante el proceso de validación técnica. & Functional Suitability & Alta \\
FR-140 & La aprobación del activo disparará la sincronización automática de los datos técnicos hacia los módulos de visualización y planificación operativa. & Functional Suitability & Alta \\
FR-141 & El sistema implementará un flujo de devolución técnica que requiera la documentación obligatoria del motivo del rechazo para fines de trazabilidad administrativa. & Functional Suitability & Alta \\
FR-142 & Tras la activación, el sistema permitirá el enriquecimiento progresivo de la ficha técnica con parámetros de rendimiento y límites operativos avanzados. & Functional Suitability & Media \\
FR-143 & El sistema generará reportes técnicos sobre el estado del alta de activos para soportar las auditorías de gestión de activos bajo ISO 55001. & Functional Suitability & Media \\
FR-149 & El sistema monitoreará los niveles de stock disponible, generando notificaciones automáticas con niveles de severidad diferenciados según la clasificación técnica del repuesto. & Functional Suitability & Alta \\
FR-150 & El tablero de gestión presentará el estado consolidado de las alertas activas, permitiendo el filtrado por prioridad, ubicación y tipo de material técnico. & Functional Suitability & Alta \\
FR-151 & El sistema integrará el flujo de alertas con la generación de requerimientos de adquisición, pre-completando la información técnica para agilizar el proceso de compra. & Functional Suitability & Alta \\
FR-152 & La resolución de condiciones de bajo stock disparará el cierre automático de las alertas, manteniendo la integridad cronológica entre la recepción y la normalización del inventario. & Functional Suitability & Alta \\
FR-153 & El sistema permitirá la consolidación periódica de requerimientos para materiales de baja rotación o criticidad menor, optimizando la gestión administrativa del almacén. & Functional Suitability & Media \\
FR-154 & El personal autorizado podrá ajustar dinámicamente los umbrales de seguridad y puntos de reposición según las variaciones estacionales o de demanda operativa. & Functional Suitability & Media \\
FR-155 & El sistema proporcionará reportes ejecutivos sobre la eficiencia del reabastecimiento, analizando los tiempos de respuesta ante alertas críticas. & Functional Suitability & Baja \\
FR-160 & El sistema permitirá la creación de nodos raíz para la taxonomía industrial, cumpliendo con los niveles jerárquicos definidos en el estándar ISO 14224. & Functional Suitability & Alta \\
FR-161 & El sistema validará que cada nueva ubicación funcional mantenga una relación jerárquica coherente con su nodo superior, impidiendo la creación de niveles desordenados. & Functional Suitability & Alta \\
FR-162 & Se proporcionarán funciones para reorganizar dinámicamente la estructura del árbol técnico, validando la integridad de las relaciones ante cada ajuste estructural. & Functional Suitability & Media \\
FR-163 & El sistema restringirá la eliminación de ubicaciones técnicas que cuenten con activos instalados o dependencias jerárquicas activas para prevenir la orfandad de datos. & Functional Suitability & Media \\
FR-164 & Se incluirán capacidades de búsqueda técnica para localizar rápidamente posiciones operativas dentro de la estructura organizacional compleja. & Functional Suitability & Baja \\
FR-165 & El sistema permitirá la generación de reportes sobre la arquitectura completa de ubicaciones funcionales para soportar procesos de planificación estratégica. & Functional Suitability & Baja \\
FR-170 & El sistema permitirá la creación de repuestos capturando su identificador único (SKU), descripción técnica, fabricante y clasificación industrial. & Functional Suitability & Alta \\
FR-171 & El sistema implementará rutinas de detección de redundancia durante el registro, analizando la similitud textual entre los repuestos nuevos y existentes. & Functional Suitability & Alta \\
FR-172 & Se proporcionará un selector jerárquico para la clasificación de materiales, permitiendo navegar a través de clases, familias y grupos técnicos. & Functional Suitability & Alta \\
FR-173 & El sistema configurará dinámicamente los campos de gestión de inventario (puntos de reorden, ubicaciones) basándose en la política de suministro del material. & Functional Suitability & Alta \\
FR-174 & El sistema permitirá la actualización de los atributos técnicos y las políticas de stock de los repuestos existentes mediante procesos debidamente auditados. & Functional Suitability & Media \\
FR-175 & Se incluirán capacidades de búsqueda avanzada para localizar materiales por múltiples criterios técnicos, optimizando la respuesta ante solicitudes masivas. & Functional Suitability & Media \\
FR-176 & El sistema permitirá generar reportes consolidados del catálogo maestro de repuestos para facilitar los procesos de inventario cíclico y auditoría. & Functional Suitability & Baja \\
NFR-121 & Las validaciones de integridad y unicidad deben realizarse en tiempo real para garantizar una interacción fluida sin bloquear la interfaz operativa. & Performance Efficiency & Alta \\
NFR-122 & El proceso de creación de activos debe ser transaccional, asegurando que ante cualquier fallo técnico se mantenga la consistencia de la base de datos maestra. & Reliability & Alta \\
NFR-123 & El sistema mantendrá la trazabilidad histórica completa de la ficha técnica, prohibiendo la eliminación física de registros para preservar la integridad de los datos de fiabilidad. & Reliability & Media \\
NFR-124 & El rendimiento del proceso de guardado y registro debe completarse en < 2s bajo condiciones estándar de conectividad industrial. & Performance Efficiency & Media \\
NFR-132 & Todas las operaciones de almacén deben procesarse de forma atómica, garantizando que los datos se mantengan íntegros ante cualquier interrupción técnica. & Reliability & Alta \\
NFR-133 & El motor de validación de existencias debe operar con latencia mínima para no retrasar el despacho de materiales al personal operativo. & Performance Efficiency & Alta \\
NFR-134 & El sistema asegurará la consistencia entre los registros de inventario y los centros de costo operativos definidos en la arquitectura del sistema. & Reliability & Alta \\
NFR-135 & El rendimiento del sistema debe permitir el procesamiento de múltiples transacciones de almacén simultáneas sin degradación del tiempo de respuesta. & Performance Efficiency & Media \\
NFR-136 & El sistema contará con mecanismos de contingencia para encolar transacciones en caso de indisponibilidad temporal de los módulos de integración. & Reliability & Media \\
NFR-144 & El sistema debe validar estrictamente la presencia de toda la documentación mandatoria antes de permitir el cambio de estado operativo. & Reliability & Alta \\
NFR-145 & El proceso de alta y activación de activos debe garantizar la integridad referencial en todos los módulos dependientes del sistema. & Reliability & Alta \\
NFR-146 & La consulta de activos pendientes de validación debe responder en < 3s incluso bajo condiciones de alta carga de datos en el sistema. & Performance Efficiency & Media \\
NFR-147 & El sistema contará con un servicio de notificaciones asincrónicas para comunicar los resultados de las revisiones técnicas al personal interesado. & Reliability & Media \\
NFR-148 & La integración de datos entre los módulos de inventario y visualización debe completarse en un tiempo que garantice la disponibilidad inmediata de la información técnica. & Reliability & Media \\
NFR-156 & La detección de condiciones críticas de stock debe realizarse con la frecuencia necesaria para garantizar la veracidad de la información operativa. & Performance Efficiency & Alta \\
NFR-157 & El sistema garantizará la entrega de notificaciones críticas a través de múltiples canales, asegurando que el personal responsable sea alertado de forma oportuna. & Reliability & Alta \\
NFR-158 & La interfaz de monitoreo de alertas debe soportar la visualización eficiente de grandes volúmenes de datos con una respuesta < 2s. & Performance Efficiency & Media \\
NFR-159 & La configuración de parámetros maestros de inventario debe realizarse de forma centralizada y entrar en vigencia inmediatamente sin interrupciones del servicio. & Usability & Media \\
NFR-166 & El motor de renderizado de la jerarquía debe soportar la visualización de estructuras extensas con una respuesta < 3s. & Performance Efficiency & Alta \\
NFR-167 & Los algoritmos de validación jerárquica deben detectar ciclos y errores estructurales en < 1s. & Reliability & Alta \\
NFR-168 & Las operaciones de reestructuración de la jerarquía deben realizarse de forma transaccional para asegurar la consistencia absoluta de los datos maestros. & Reliability & Media \\
NFR-169 & El sistema implementará mecanismos de control de concurrencia para evitar conflictos ante la edición simultánea de la estructura técnica por múltiples ingenieros. & Reliability & Media \\
NFR-177 & La validación de unicidad de los identificadores de repuestos debe realizarse en < 1s para asegurar la fluidez del proceso administrativo. & Performance Efficiency & Alta \\
NFR-178 & El motor de detección de duplicados debe operar con un tiempo de respuesta < 2s para minimizar la creación de registros redundantes. & Performance Efficiency & Alta \\
NFR-179 & La validación taxonómica debe realizarse en < 1s, empleando mecanismos de almacenamiento temporal para optimizar la consulta de las tablas maestras de categorías. & Performance Efficiency & Alta \\
NFR-180 & El sistema debe garantizar un rendimiento de búsqueda constante incluso ante catálogos de materiales extensos. & Performance Efficiency & Media \\
NFR-181 & El sistema asegurará la compatibilidad y conversión de unidades de medida para mantener la consistencia técnica en todas las transacciones de inventario. & Compatibility & Media \\
\bottomrule
\end{longtable}
\normalsize